(1) 碰撞期间悬线的约束冲量沿竖直方向，极短时间内的水平外冲量可以忽略。于是水平方向动量守恒 \(mv=(M+m)V\)，得到嵌入后共同水平速度 \(V=mv/(M+m)\)；这一等式不要求碰撞前后机械能相等。
(2) 碰撞结束后将弹丸与摆球视为同一质量 \(M+m\) 的摆体，沿无摩擦轨道上升时机械能守恒。最高点速度为零，质心抬升 \(h=L(1-\cos\theta)\)，所以 \((M+m)V^2/2=(M+m)gL(1-\cos\theta)\)，即 \(V=\sqrt{2gL(1-\cos\theta)}\)。
(3) 联立得到 \(v=(M+m)\sqrt{2gL(1-\cos\theta)}/m\)。完全非弹性碰撞期间损失的动能为 \(\Delta E=mv^2/2-(M+m)V^2/2=mMv^2/[2(M+m)]>0\)，对应不可逆形变、声和内能增加，与后续摆动守恒不能混淆。
(4) 若 \(m\ll M\)，有 \(V\simeq(m/M)v\)；当 \(\theta\ll1\) 时 \(1-\cos\theta\simeq\theta^2/2\)，因此 \(v\simeq(M+m)\sqrt{gL}\theta/m\)。碰撞时只使用水平冲量平衡，碰后再使用机械能守恒，这种阶段划分是求解的关键。
